\section{Teoría de diseño y cálculos}

\subsection{Pertenencia a un polígono}

El motor decide si un punto está dentro de un polígono con el método de lanzamiento de rayos
(\emph{ray casting}) y la regla par-impar: se traza un rayo desde el punto y se cuenta cuántos
lados del polígono cruza; si el número es impar, el punto está dentro. Los vértices se dan en
orden y el último se une con el primero. En un polígono que se corta a sí mismo, la región
cubierta dos veces cuenta como fuera.

Un punto que cae exactamente sobre el borde se clasifica con una regla semiabierta: en un
rectángulo alineado con los ejes, el lado sur, el lado oeste y la esquina suroeste cuentan como
dentro, y el resto, como fuera. El resultado sobre el borde exacto no debe usarse como dato.

\subsection{Distancia al borde}

La distancia que se muestra es la distancia al lado más cercano, no al vértice más cercano. Para
calcularla, el motor proyecta las coordenadas a un plano local centrado en el punto:

\begin{equation}
x = (\lambda - \lambda_0)\,\frac{\pi}{180}\,R\cos\varphi_0, \qquad
y = (\varphi - \varphi_0)\,\frac{\pi}{180}\,R,
\end{equation}

con $R = 6371$~km, de modo que un grado de latitud equivale a unos 111\,195~m. La distancia es
el mínimo, entre todos los lados, de la distancia del origen al segmento. Es una distancia sin
signo: la pertenencia al polígono se decide aparte.

El error de esta proyección plana respecto a la distancia de haversine está documentado en el
motor como menor que 0.2~\% hasta los 85° de latitud, y crece hacia los polos. No se usa cerca de
los polos, y tampoco con polígonos que crucen el antimeridiano (longitud $\pm180^\circ$), porque
las longitudes se restan como números simples.

\subsection{Una sola llamada que no confunde errores con ``fuera''}

Las funciones de pertenencia devuelven un valor lógico, así que no pueden distinguir ``fuera'' de
``entrada inválida''. Por eso la aplicación usa \texttt{geofence\_check()}, que en una sola llamada
devuelve si el punto está dentro y la distancia al borde, y devuelve un código de error si el fix
no es válido (coordenadas fuera de rango o no numéricas, polígono vacío o motor sin iniciar). Un
fix inválido nunca se interpreta como una salida de la cerca.

\subsection{Cerca cuadrada alrededor de un punto}

Para un centro $(\varphi_0, \lambda_0)$ y una mitad de lado $h$, las esquinas se calculan con

\begin{equation}
\Delta\varphi = \frac{h}{111\,195}, \qquad \Delta\lambda = \frac{\Delta\varphi}{\cos\varphi_0},
\end{equation}

y las esquinas son $(\varphi_0 \mp \Delta\varphi,\ \lambda_0 \mp \Delta\lambda)$, en el orden
suroeste, sureste, noreste y noroeste. La longitud se ajusta con el coseno de la latitud porque un
grado de longitud mide menos al alejarse del ecuador. La construcción se rechaza si la latitud del
centro pasa de 85° o si el cuadrado cruzaría el antimeridiano.

\subsection{De la distancia al color}

Con una distancia segura $S = 20$~m y una distancia $d$ al borde, con el punto dentro, el color es

\begin{equation}
t = \min\!\left(\frac{d}{S},\,1\right), \qquad
R = \operatorname{round}\bigl(255\,(1-t)\bigr), \qquad
G = \operatorname{round}\bigl(255\,t\bigr), \qquad B = 0.
\end{equation}

A $d \ge S$ el color es verde puro, y a $d = 0$, rojo fijo. Fuera de la cerca, o ante cualquier
entrada inválida (distancia no finita o negativa, o $S \le 0$), el resultado es rojo
parpadeante, el estado a prueba de fallos.

\subsection{Presupuesto de tiempo}

El costo de \texttt{geofence\_check()} es lineal en el número $N$ de vértices, porque recorre cada
lado una vez para decidir la pertenencia y otra para calcular la distancia. Medido en la placa
(ver la sección de resultados), el tiempo por llamada se ajusta a

\begin{equation}
T(N) \approx \dBenchFitOffsetUs~\mu\text{s} + \dBenchFitSlopeUs~\mu\text{s}\cdot N,
\end{equation}

con lo que el límite de 50~ms de SR-12 se alcanzaría cerca de los \dBenchLimitVertices{} vértices.
El motor usa números de doble precisión; como la unidad de coma flotante del Cortex-M33 es de
precisión simple, es probable que esas operaciones se ejecuten por software y que ahí esté el
costo por vértice. Esto es una inferencia a partir de los indicadores del compilador, no una
medición.

\subsection{Resolución del reloj de medida}

El contador de ciclos del kernel corre a \dClockHz~Hz, de modo que cada lectura tiene una
resolución de \dTickUs~$\mu$s. Una llamada puede durar menos que eso, por lo que no se mide una
llamada suelta sino un lote de 40 y se divide el tiempo total entre 40.
